\section{Bootstrapping basic constructs}
\label{bootstrapping-basic}

\begin{note}
can this section be shortened further?
\end{note}

The base syntax of ILD is insufficient for writing nontrivial programs
directly. However, its \(\ildsf{macroexpand}\) facility allows
extending the syntax with arbitrary metaprogramming constructs, which can
be built incrementally.
In this section, we present a standard library
of such constructs, culminating in \ild{letrec}.
Keeping that complexity out of the core language makes it simple and portable
across implementations.

\subsection{Lambda}
\label{lambda-macro}
First, we define a macro called \ild{lambda} that will let us introduce new
user-defined functions (\cref{abstraction}). We want expansion of the \ild{lambda}
macro to result in a call to \ild{mk-lambda} given the current binding environment
(\ild{free-vars}) and the quoted argument names and body (\cref{lambda-expansion-example}):
\begin{listingbox}\begin{ildcode}
(!lambda (x y) (+ x y))
; should become
(mk-lambda (free-vars) (quote (x y))
  (quote (+ x y)))
\end{ildcode}\end{listingbox}
\lstcaption{Example expansion for \ild{!lambda} expression}\label{lambda-expansion-example}
The outer call to \ild{mk-lambda} introduces a new \(\ildsetUserFunc\)-object
carrying the current binding environment
and a body. The body then, when called, in turn
expands to an expression embedding \ildhost{mk-lambda} in its head.
\begin{listingbox}\begin{ildcode}
(mk-lambda
  (free-vars)
  '(arg-names body)
  '(cons mk-lambda
    (cons (cons free-vars '())
      (cons (cons quote (cons arg-names '()))
        (cons (cons quote (cons body '()))
          '())))))
\end{ildcode}\end{listingbox}
\lstcaption{Implementation of the lambda macro}\label{lambda-macro-def}

\subsection{Basic utilities}
We define helper utilities in a scope where the \ild{lambda} macro is present.
\begin{listingbox}\begin{ildcode}
(!lambda (p) (car (cdr p))) ; body of 'cadr'
(!lambda (p) (car (car p))) ; body of 'caar'

; 'Y' is special case of poly-fix
(!lambda (f) (car (poly-fix f)))
\end{ildcode}\end{listingbox}
\lstcaption{Definitions of basic helper primitives}
Since \ild{lambda} is an ordinary function, it can also be called directly
to obtain the expression it expands to, which is useful when writing macros
that generate \ild{(!lambda <args> <body>)} expressions.
We can now implement branching via the builtin function
\ildhost{bool-to-k}.
To avoid evaluating both branches, we delay the evaluation of
each branch until the condition is known by wrapping it in a 0-argument
function, and force the selected branch afterwards:
\begin{listingbox}\begin{ildcode}
; body of 'if'
(!lambda (cond t f)
  ; generate an expression that
  (list     ; forces the result
    (list   ; applies the selector,
            ; obtaining a delayed value
      (list bool-to-k cond) ; selector

      ; delayed #t and #f branches
      (lambda '() t)
      (lambda '() f))))
\end{ildcode}\end{listingbox}
\lstcaption{Definition of the \ild{if} macro}\label{if}
Since \ild{letrec} is not available at this point in the bootstrap process,
recursive functions such as \ild{map} can be defined via \ild{Y}:
\begin{listingbox}\begin{ildcode}
; body of 'map'
(Y (!lambda (map)
  (!lambda (f l)
    (!if (null? l)
      l
      (!if (pair? l)
        (cons
          (f (car l)) (map f (cdr l)))
        (make-fail
          (list 'not-a-list l)))))))
\end{ildcode}\end{listingbox}
\lstcaption{Definition of the \ild{map} function}\label{def-map}

\subsection{Letrec}
\label{letrec}
With these building blocks we can define a
macro called \ild{letrec} that transforms a list of (possibly mutually-recursive)
function definitions into a call to \(\ildsym{poly-fix}\):
\begin{listingbox}\begin{ildcode}
(!letrec (
    (even? (!lambda (n)
      (!if (= n 0) #t (odd? (- n 1)))))
    (odd? (!lambda (n)
      (!if (= n 0) #f (even? (- n 1)))))
    (main (!lambda () (even? 42))))
  (main))
; becomes
(apply
  (!lambda (even? odd? main) (main))
  (poly-fix
    (!lambda (even? odd? main)
      (!if (= n 0) #t (odd? (- n 1))))
    (!lambda (even? odd? main)
      (!if (= n 0) #f (even? (- n 1))))
    (!lambda (even? odd? main) (even? 42))))
\end{ildcode}\end{listingbox}
\lstcaption{A \ild{letrec} example and its expansion into a \ild{poly-fix} application}\label{letrec-expansion}
The implementation of \ild{letrec} is merely technical and is omitted.
